\chapter{Introduction}
\label{chap:introduction}


\section{Inkwell Data Ltd}
This thesis project was carried out at \textit{Inkwell Data Ltd}, a company
% add 6 months stage for MSc Computer Science
operating in the \textit{Industrial Internet of Things} (\textit{IIoT}) sector,
specializing in connectivity solutions and digital-twin technologies for
industrial assets. In this context, a digital twin is understood as a dynamic
digital representation of a physical device and its software representation,
continuously connected to the corresponding real-world object for data exchange
thoughout its lifecycle.  
Founded in 2019, with teams interconnecting London and Padua, the company
started with \textit{Altior}, a middleware \textit{Platform as a Service}
(\textit{PaaS}) designed to simplify the creation and management of digital
twins at an enterprise scale.  

\begin{figure}[H]
      \centering
      \includegraphics[alt={Logo \textit{Inkwell Data Ltd.}}, width=0.25\columnwidth]{img/inkwell_data_logo.jpeg}
      \caption{\textit{Inkwell Data Ltd} logo}
      \label{fig:inkwelldata}
\end{figure}

At the core of \textit{Altior}, device twin instances are treated as data
structures reflecting the current state of the corresponding real-world object,
allowing the platform to remain highly responsive even under heavy load. Rather
than focusing on the mere collection of data, \textit{Altior} places particular
emphasis on how this data flows and is processed once collected, for instance by
publishing device updates onto message queues for downstream consumption. This
% avoid use of instance and provide more exaples
focus on data flow is further supported by \textit{Altior}'s distributed nature:
% remove further supported, just say distributed logic
the platform can be transparently replicated and clustered across multiple
% remove transparently replicated, add horizontal and vertical section of white papers
interconnected servers, providing horizontal scalability, load balancing and
fault tolerance, all while maintaining a strong emphasis on security.

On top of the Altior platform, Inkwell Data develops \textit{Software as a
Service} (\textit{SaaS}) applications tailored to specific industrial
domains. One such application is \textit{Swarm}, a product built on Altior
and primarily focused on the water metering sector, although the company's
technology is designed to generalize to a broader range of connected
industrial devices.

% add swarm organizations in elixir apps


\section{Context and Motivation}
Over the last decades, the way in which software has been thought about has go  

\section{Problem Statement and Objectives}

\section{Thesis Structure}
% Completely to be remade
This thesis is organized into seven chapters, covering the background,
requirements, architecture, implementation, and evaluation of DEX.

\begin{description}
      \item[\textbf{Chapter 2}] provides the background required to understand
      the project. It introduces data exploration, distributed data systems,
      heterogeneous data sources, data abstraction and governance, and
      distributed query processing. The chapter concludes with an overview of
      existing solutions and technologies related to the presented problem.
      \item[\textbf{Chapter 3}] describes the requirements and design goals of
      DEX. It defines the functional requirements, quality requirements, and
      constraints that guided the development of the system.
      \item[\textbf{Chapter 4}] presents the architecture of DEX. It first
      provides an overview of the system and its main architectural principles,
      and then describes its adoption of the Hexagonal Architecture. The chapter
      also discusses distributed execution, fault tolerance, and the application
      of microservice principles within the BEAM environment.
      \item[\textbf{Chapter 5}] describes the implementation of DEX. It
      introduces the main technologies adopted by the system and details the
      implementation of query processing, data source integration, persistence,
      configuration, and client interaction. The chapter also discusses the
      concrete mechanisms used to provide fault tolerance.
      \item[\textbf{Chapter 6}] evaluates the developed system. It describes the
      experimental setup and presents the results of the performance and
      scalability evaluations. The obtained results are subsequently discussed
      to assess the behavior of DEX under the considered experimental
      conditions.
      \item[\textbf{Chapter 7}] concludes the thesis by summarizing the main
      contributions of the project and discussing its current limitations. The
      chapter concludes by presenting possible directions for future
      development.

% In merito alla stesura del testo, all'interno del presente documento sono state adottate le seguenti convenzioni tipografiche:
% \begin{itemize}
%       \item gli acronimi, le abbreviazioni e i termini ambigui o di uso non comune menzionati vengono definiti nel glossario;
%       \item per la prima occorrenza dei termini riportati nel glossario viene utilizzata la seguente nomenclatura: \textit{parola}\glox\gloxspacing;
%       \item i termini facenti parti del gergo tecnico sono evidenziati con il carattere \textit{corsivo}.
% \end{itemize}
\end{description}

% \section{Organizzazione del testo}
% Nel presente capitolo viene presentata l'introduzione della tesi, fornendo una panoramica sull'azienda, sul contesto applicativo,
% sul progetto, sugli strumenti utilizzati e sul prodotto portato a termine durante lo svolgimento dello stage. \newline
% In seguito il documento presenterà la seguente organizzazione:

% \begin{description}
%       \item Il {\hyperref[chap:analisi-requisiti]{Capitolo 2}} descrive la fase di analisi dei requisiti che è stata
%             svolta dall'azienda in fase antecedente all'inizio dell'attività di stage, in modo da permettere una comprensione più profonda
%             di quelle che sono le necessità da soddisfare e gli obiettivi da raggiungere all'interno di questo progetto;

%       \item Il {\hyperref[chap:progettazione]{Capitolo 3}} illustra l'attività di progettazione che è stata svolta in vista dell'implementazione
%             dei grafici prodotti, definendo e individuando le soluzioni implementative che sono state attuate durante la successiva attività di codifica;

%       \item Il {\hyperref[chap:realizzazione-testing]{Capitolo 4}} approfondisce l'attività di codifica delle componenti grafiche presentate
%             all'interno della libreria implementata e dei relativi test effettuati per garantire la qualità e la funzionalità del prodotto finale;

%       \item Il {\hyperref[chap:rilascio]{Capitolo 5}} descrive l'attività di rilascio dell'SDK realizzato, presentando le modalità di integrazione
%             all'interno di un prodotto \textit{Datasoil S.r.l.};

%       \item Il {\hyperref[chap:conclusioni]{Capitolo 6}} presenta un epilogo del progetto svolto, includendo un consuntivo finale delle attività svolte,
%             una valutazione del progetto e una sezione dedicata ai possibili sviluppi futuri del prodotto.
%             La sezione infine si conclude con una riflessione finale sull'esperienza di stage svolta.
% \end{description}

